(1) 二维波矢态数每单位面积为 \(g\,d^2k/(2\pi)^2\)。热激发单位面积能量 \(u=g/(2\pi)\int_0^\infty k\,dk\,\hbar ck/(e^{\beta\hbar ck}-1)\)。令 \(x=\beta\hbar ck\)，并用 \(\int_0^\infty x^2/(e^x-1)dx=2\zeta(3)\)，得 \(u=g\zeta(3)(k_BT)^3/[\pi\hbar^2c^2]\)。紫外端由 Bose 因子指数收敛；零温真空能已按题设扣除。

(2) 线性色散在二维使无质量各向同性激发的应力迹满足 \(2P=u\)，因此 \(P=u/2=g\zeta(3)(k_BT)^3/(2\pi\hbar^2c^2)\)。也可用巨势按空间尺寸的齐次缩放验证。此处 \(P\) 是二维力每长度，其乘以面积具有能量量纲；相应焓 \(H=U+PA=3U/2\)。不要把三维辐射压 \(P=u/3\) 直接用于二维系统。

(3) 在化学势零且无其他守恒粒子数时，热力学关系 \(U=TS-PA\) 给 \(S=(U+PA)/T=3U/(2T)=3g\zeta(3)Ak_B^3T^2/(2\pi\hbar^2c^2)\)。在固定面积下 \(C_A=(\partial U/\partial T)_A=3U/T\)，热容随 \(AT^2\) 缩放，与二维线性色散的态密度成正比。熵表达式具有单位 J/K，符合广延性。

(4) 准静态绝热演化保持 \(S\propto AT^2\) 不变，故 \(T\sqrt A=\mathrm{const}\)。等价地，\(P\propto T^3\)，结合绝热线得到 \(PA^{3/2}=\mathrm{const}\)。三维光子气体则满足 \(S\propto VT^3\)、\(TV^{1/3}=\mathrm{const}\)；维数不同导致指数不同。若存在质量间隙、固定粒子数或非线性色散，以上二维幂律将不再普适。
